\chapter{Impacts and Design Constraints of the Project}
\label{ch:impacts}

This chapter examines the non-technical constraints that shaped LectureAssist,
economic, social, legal, ethical and environmental, and the impact the system
is likely to have once deployed. Several of the most consequential engineering
decisions in this project, notably the refusal to depend on any paid cloud
API, were driven by these constraints rather than by technical considerations.
They are therefore not an afterthought appended to a finished design; they are
among the inputs that produced it.

\section{Economic Impact}
\label{sec:impacts-economic}

\paragraph{The constraint that drove the architecture.} The system was
required to run without a paid API. This is not a stylistic preference, and it is
worth being precise about why. Per-token cloud pricing makes routine, high-volume
practice, which is precisely the usage pattern that makes retrieval practice
effective, prohibitively expensive for students in low- and middle-income
settings. It also makes the tool's availability contingent on a credit card and an
international payment rail that many students in Bangladesh and comparable settings
simply do not have; the barrier is not only the price but the mechanism of payment.
Every model in LectureAssist is therefore open-weight and locally served. The
consequence is a hard VRAM budget, which is why the generator is a $4$B model rather
than a frontier one, and much of the engineering documented in
Chapter~\ref{ch:methodology} exists to extract reliable structured output from a
small model. The evaluation in Chapter~\ref{ch:results} can be read as a test of
whether that constraint is survivable: it finds that a well-structured agentic loop
around a $4$B model measurably outperforms a single call to the same model, which is
the mechanism by which the cost constraint is absorbed rather than merely accepted.

\paragraph{Where the cost actually goes.} Removing the per-token bill does
not make the system free; it moves the cost from a recurring charge to a
one-off hardware requirement, and it is worth being clear about the trade.
A cloud deployment charges per question generated, so cost scales with
exactly the behaviour the system is trying to encourage: a student who
practises more pays more. Local inference inverts that. The marginal cost of
the thousandth question is the electricity to generate it, which means the
usage pattern that makes retrieval practice effective is the pattern the cost
model rewards. What replaces the bill is a floor: the system needs a GPU with
enough memory to hold a small generator and a larger judge, which in this
project is met by the free tier of a hosted notebook rather than by hardware
the student owns.

\paragraph{The compute cost of the agentic loop.} The architecture is not free
of economic consequence either. As Chapter~\ref{ch:results} records, the
five-agent pipeline costs roughly an order of magnitude more compute than the
single call it outperforms, because planning, retrieval, validation and retry
each add model invocations. On a paid API that multiplication is what makes
multi-agent decomposition expensive, and it is why the surveyed systems that
adopt it also depend on commercial inference. On local inference the same
multiplication costs time rather than money, which is the specific reason the
design is affordable here and would not be affordable in the cloud. The
ablation designs of Section~\ref{sec:res-ablation} exist partly to establish
whether the whole multiplication is necessary or whether a cheaper subset of
the agents captures most of the benefit, which is an economic question as much
as a scientific one.

\paragraph{Institutional adoption.} For a department rather than an individual,
the relevant comparison is against the staff time currently spent writing
practice items by hand. The system does not replace that work, because its
output is unreviewed until a human reviews it, but it changes the task from
authoring to checking. Whether that is a saving depends on how much of the
generated set survives review, and this project measures acceptance under an
automated judge rather than under an instructor, so the honest position is
that the institutional case is plausible but not evidenced here.

\section{Impact on Societal, Health, Safety, Legal and Cultural Issues}
\label{sec:impacts-society}

\paragraph{Educational and societal impact.} The core societal argument for
LectureAssist is one of access. Retrieval practice, repeatedly testing oneself on
material rather than re-reading it, is among the most effective study strategies
known, and the effect is large, replicated and long
established~\cite{roediger2006testenhanced,karpicke2011retrieval}. In the most
comprehensive review of study techniques available, practice testing is one of only
two techniques rated as having high utility across learners, materials and
criterion tasks~\cite{dunlosky2013improving}. Yet a supply of good practice questions
is a resource distributed very unevenly. A student at a well-resourced institution
has question banks, teaching assistants and instructor office hours; a student
without those has the end-of-chapter exercises and nothing else, and those exercises
are not tied to what their particular instructor actually emphasised in class. By
generating unlimited, validated, difficulty-controlled questions from the student's
\emph{own} lecture material, the system puts a form of practice that was previously a
privilege of well-resourced institutions within reach of anyone with a laptop. This
aligns directly with UN Sustainable Development Goal~4 (Quality Education) and,
because it removes a cost barrier rather than an ability barrier, with Goal~10
(Reduced Inequalities).

\paragraph{Safety: the risk of confidently teaching the wrong thing.} The most
serious safety issue in a system of this kind is not a crash: it is a fluent,
well-formed multiple-choice question whose marked answer is \emph{wrong}. A student
who trusts the tool will learn the error, rehearse it through exactly the retrieval
practice that makes the tool valuable, and carry it into an exam. The tool's fluency
works actively against their catching it, because the surface features a student uses
to judge whether a question is trustworthy, clear wording, plausible distractors,
a confident explanation, are precisely the features a language model produces most
reliably, whether or not the marked option is correct.

This risk shaped the evaluation methodology rather than merely being noted in it.
Answer correctness is measured separately from question quality, by an Answer
Generator agent that re-solves each question without ever seeing the marked option
(Section~\ref{sec:meth-answer}), and the solver is itself calibrated against
questions whose answers are known independently, so that its verdict is not
over-trusted. The honest conclusion is that the system's output is suitable for
\emph{practice} and must not be treated as an authoritative answer key. Any
deployment should surface this to the student directly, and no generated question
should be used for graded assessment without human review. The design implication is
concrete: the interface shows the source timestamp for every question, so that a
student who doubts an answer can jump to the moment in the lecture it came from
rather than having to trust or discard the question wholesale.

\paragraph{Legal and privacy considerations.} Lecture recordings and course materials
are frequently copyrighted by the instructor or the institution, and lecture audio
contains the identifiable voices of the instructor and, incidentally, of students who
ask questions. Uploading such material to a third-party cloud service to be processed, and potentially retained, logged, or used for training, raises copyright and
data-protection questions that most students are not in a position to evaluate, and
that they are typically not even prompted to consider. Because LectureAssist runs
entirely locally, the lecture never leaves the machine the user controls: there is no
third-party processor, and therefore no third-party retention. Users must still hold
the right to use the material they upload, and the system should not be used to
redistribute copyrighted lectures. The persistent per-topic performance profile is
educational data about an identifiable student, and it is stored in the user's own
storage rather than on any server operated by the developers.

One caveat must be recorded honestly. The evaluation deployment exposes the backend
through a public tunnel without authentication, as noted in
Section~\ref{sec:concl-limits}. That is acceptable for a supervised demonstration and
is not acceptable for real use, and it is the first thing that would have to change
before the system were put in front of students.

\paragraph{Cultural and linguistic considerations.} OCR is configured for both
English and Bengali, reflecting the reality of instruction in Bangladesh, where
lecturers routinely switch languages mid-sentence and write on the board in either.
A pipeline that only read English would silently discard part of the lecture, and
would do so invisibly, since a missing board block looks exactly like a board that
was never written on. This is a recurring theme in the system's failure modes: the
dangerous errors are the ones that produce no error, and the design response is to
widen coverage rather than to detect absence, because absence in this pipeline is
undetectable in principle.

\section{Ethical Considerations}
\label{sec:impacts-ethics}

Ethical questions arise for this system in two distinct registers, and it is
worth separating them. The first concerns the conduct of the research itself:
what this report claims, and on what evidence. The second concerns the artifact:
what a deployed question generator can do to the people who rely on it. The
safety and privacy dimensions of the second register are treated in
Section~\ref{sec:impacts-society}; this section addresses the commitments and
obligations that follow from them.

\subsection{Research Ethics}

Two ethical commitments governed this work, and both are visible in the
structure of Chapter~\ref{ch:results}.

\paragraph{No fabricated results.} Every number reported in
Chapter~\ref{ch:results} comes from a measured run, and the report claims
nothing beyond what was measured. It would have been easy to present a fuller
looking set of results by supplying plausible values for comparisons that were
designed but not carried out; that option was rejected, and the chapter reports
the comparison that was actually run and analyses it thoroughly instead. The
limitations of those measurements, particularly the saturation of the judged
dimensions and the weakness of an LLM solver as a proxy for correctness, are
stated plainly rather than omitted because they are unflattering. Where a labelling correction had to be applied
that cannot be re-derived from the stored artifacts, that fact is disclosed in
the results chapter and again in the limitations, rather than quietly
incorporated.

\paragraph{No misrepresentation of the system's capability.} The tool is an
aid to practice, not a substitute for the instructor or the textbook, and this
report is explicit about where its output cannot be trusted. A system that
generates assessment material carries a specific professional obligation not
to overstate its reliability, because the people most likely to over-trust it
are precisely the students with the least access to an expert who could
correct them: the same students the system is built for.

\subsection{Ethics of Deployment}

\paragraph{Academic integrity.} A system that turns a lecture into questions can
also be turned toward producing answers, and the boundary between the two is
thinner than it looks. LectureAssist is deliberately built as a
\emph{self-assessment} instrument: it generates questions for a student to
attempt, grades the attempt, and reports which topics were weak. It does not
draft submittable coursework, and the short-answer grader returns feedback
against the lecture rather than a polished answer to be copied. That boundary is
a design decision rather than a technical impossibility, and it is stated here so
that it is understood as a commitment which any future extension of the system
inherits. An instructor adopting the tool should also be told what it is: an
item generator whose output is unreviewed until a human reviews it.

\paragraph{The obligation created by confident error.} The safety analysis in
Section~\ref{sec:impacts-society} establishes that a fluent, wrong question is
the characteristic failure of this class of system. The ethical consequence is
that the burden of detection cannot be placed on the learner, because a student
who could reliably detect a subtly wrong question about a topic would not need
to be tested on it. This is why answer correctness is checked by a second agent
that never sees the marked option, why items that fail validation are flagged as
overrides rather than silently accepted, and why every generated item carries a
timestamp back to the source so a doubt can be resolved against the lecture
rather than against the model. These are ethical requirements expressed as
architecture, and they are the reason for design choices that would otherwise
look like unnecessary cost.

\paragraph{Data, consent and ownership.} Lecture recordings contain the voice,
and often the image, of an instructor who has not consented to third-party
processing, and the material itself is usually the instructor's or the
institution's intellectual property. Running entirely on local inference is the
strongest available mitigation, because the recording is never transmitted to a
provider who could retain it, train on it, or be compelled to disclose it. The
legal dimension is treated in Section~\ref{sec:impacts-society}; the ethical
point is narrower and applies even where the law is silent: a student's right to
study a lecture they attended does not extend to redistributing their
instructor's teaching, and the system should not make redistribution the path of
least resistance. Retention is correspondingly limited to what a session needs.

\paragraph{Bias, coverage and who the system serves.} Two kinds of bias bear on
this system. The first is inherited: the generator, the judge and the
independent solver are all language models trained on corpora that
over-represent English and well-resourced academic domains, so question quality
will not be uniform across subjects or languages, and this report offers no
evidence about performance outside the single English STEM chapter evaluated in
Chapter~\ref{ch:results}. The second is structural and specific to the
evaluation: the judge and the solver are drawn from the same model family as the
generator, so they share its blind spots, and agreement between them is evidence
of consistency rather than of correctness. That limitation is stated in
Chapter~\ref{ch:results} as a measurement caveat; it is repeated here because it
is also an ethical one. A system that grades itself with a model like itself
will systematically fail to detect the errors that model is prone to, and it
should not be presented to students as an independent check on their
understanding.

\paragraph{Boundaries of responsible use.} Taken together, the commitments above
define what this system may honestly be claimed to do. It is a practice and
revision aid, grounded in a specific lecture, whose output is checked but not
guaranteed. It is not an examination instrument, not a substitute for the
instructor, and not a source of authority about the subject matter. Deploying it
in any of those three roles would exceed what the evaluation in
Chapter~\ref{ch:results} supports, and the project treats that as an ethical
limit rather than a marketing one.

\section{Impact on Environment and Sustainability}
\label{sec:impacts-env}

The environmental footprint of an AI system is dominated by the energy consumed
during inference at scale. LectureAssist's design happens to be favourable on this
axis, largely as a side effect of the constraint that it must run on a free-tier GPU.

\paragraph{Small models, local inference.} Question generation runs on a
$4$B-parameter model rather than a frontier-scale one, and the entire system runs on
a single consumer GPU. Inference on a small local model consumes orders of magnitude
less energy per request than routing that request to a hyperscale data centre running
a model two to three orders of magnitude larger, and it avoids the network transfer
and data-centre cooling overhead entirely. No model was trained or fine-tuned during
this project, all models are used as released, so the substantial carbon cost
of training was not incurred by this work.

\paragraph{Caching and resumability as energy savings.} OCR results are cached and
keyed by a hash of the video, so re-processing the same lecture costs nothing.
Consecutive frames whose extracted text is more than $78\%$ similar are treated as
one board state, so a slide left on screen for two minutes contributes one block
rather than sixty. Bulk evaluation runs are checkpointed to persistent storage so
that a disconnection resumes rather than restarts. Each of these was implemented for
practical reasons, free-tier compute is scarce and Colab runtimes die, but each
also directly avoids recomputing work that has already been done, which is the
cheapest energy saving available and the only one that costs nothing in quality.

\paragraph{The honest counterweight: the agentic loop is not free.} The agentic
pipeline consumes substantially more compute than a single-shot baseline, because it
issues a planning call, up to three drafting calls, a validation call per attempt,
and any retrieval calls the generator elects to make, against the baseline's one.
Chapter~\ref{ch:results} finds that this buys a real and consistent quality
improvement on every metric that discriminates, so the expenditure is not wasted.
But ``not wasted'' is a weaker claim than ``efficient'', and the report should not
pretend otherwise. This is exactly why every ablation and merging table in
Chapter~\ref{ch:results} carries a mandatory cost column: a configuration that
matches the full pipeline's quality using two-thirds of the calls is
environmentally as well as economically preferable, and the experiment is designed so
that such a configuration would be identified rather than overlooked. Reducing the
call count at fixed quality is simultaneously a performance goal, a cost goal and an
environmental one, and the three do not conflict.

\paragraph{Sustainability of access.} A tool that requires a subscription is
sustainable only for as long as the student can pay for it, and only for as long as
the vendor chooses to keep that model available at a stable price. Models are
deprecated, prices change, and free tiers are withdrawn. Because LectureAssist
depends on no external service, it will continue to work on hardware a student
already owns, for as long as they own it, with weights they have already downloaded.
This is an important form of sustainability that is easy to overlook when
sustainability is framed purely in terms of energy consumption: a system that stops
working when a company changes its pricing has a lifespan measured in business
decisions rather than in hardware.
